\section{Introduction}
\label{sec:introduction}

Physics-informed neural networks (PINNs) approximate solutions of partial
differential equations (PDEs) by minimizing losses that combine data mismatch,
boundary-condition error, and the differential-equation residual evaluated at
collocation points \cite{raissi2019physics,karniadakis2021physics}.  Because the
residual term is an empirical approximation to a domain integral, the placement
of collocation points directly changes the optimization problem seen by the
network.  This dependence becomes especially important on domains with
perforations, obstacles, narrow regions, or localized solution features, where a
fixed uniform point set may allocate too little resolution to the parts of the
domain that dominate the error or residual.

Adaptive collocation strategies address this issue by changing the residual
point distribution during training.  A systematic comparison of PINN sampling
strategies, including random, low-discrepancy, and residual-adaptive methods, is
given by \cite{wu2023comprehensive}.  Residual-based adaptive distribution (RAD)
methods increase the sampling probability of candidate points with large PDE
residuals, while low-discrepancy sequences such as Halton and Sobol provide
strong non-adaptive coverage baselines
\cite{halton1960efficiency,sobol1967distribution}.  These approaches represent
different tradeoffs.  Pointwise residual weighting can respond quickly to
localized defects, but it may also produce highly concentrated sampling
distributions when the residual field contains sharp outliers.  Low-discrepancy
sampling provides stable global coverage, but it does not use information from
the current model.  For complex geometries, a useful adaptive sampler should
therefore combine residual sensitivity with an explicit mechanism for preserving
spatial coverage.

This paper introduces a fixed-budget, mesh-scaffolded adaptive collocation
method for PINNs on complex geometries.  The method evaluates the current PINN
residual on a simplicial mesh, smooths and rank-normalizes elementwise residual
scores, and samples a fixed number of interior collocation points from a
power-tempered element distribution.  The mesh is used as a scaffold for
geometry-aware residual aggregation and optional refinement; it is not used to
increase the training budget.  This separation is important experimentally:
adaptive sampling is compared against random, RAD, and low-discrepancy baselines
under matched collocation budgets and matched resampling schedules.

We evaluate the method in \pinnpoint, a controlled benchmark framework that
couples finite-element reference solutions with matched-budget PINN training.
All methods in a comparison share the same supervised FEM data, reference
solution, network architecture, training budget, and seed set; only the
collocation policy changes.  The experiments cover Poisson, advection-diffusion,
Navier-Stokes, and Allen-Cahn benchmarks.  Rather than treating adaptive
sampling as uniformly superior, we report the empirical tradeoffs among solution
error, fixed-reference residual, and runtime, including seed-paired statistical
tests for the main comparisons.

The main contributions of this work are:
\begin{itemize}
    \item A fixed-budget adaptive collocation algorithm that uses a simplicial mesh scaffold to aggregate residual information while preserving explicit control over the number of residual points.
    \item A power-tempered, rank-persistent element distribution that balances residual concentration against geometric coverage.
    \item A matched-budget evaluation protocol for comparing adaptive, random, RAD, and low-discrepancy collocation policies under common FEM supervision and reference metrics.
    \item A multi-seed empirical study across four PDE benchmarks, reporting solution error, fixed-reference residual, runtime, and paired statistical comparisons.
\end{itemize}

The remainder of this paper is organized as follows.  \Cref{sec:method}
defines the PINN objective, the mesh-scaffolded sampling procedure, and the
comparison samplers.  \Cref{sec:experiments} describes the benchmark problems
and experimental protocol.  \Cref{sec:results} presents the empirical results
and tradeoff analysis, and \cref{sec:conclusions} summarizes the conclusions.
